%	JASA LaTeX Sample File, Preprint Sample
%
%  Beginner Latex users should refer to their favorite online documentation
%  here is one from the TeX Users Group 
%	https://www.tug.org/twg/mactex/tutorials/ltxprimer-1.0.pdf
%
%  Useful FAQ from  https://journals.aps.org/revtex/revtex-faq
% 

%%%%%%% For Preprint
%% For manuscript, 12pt, one column style

%  \documentclass[preprint]{JASAnew}

%%%%% Preprint Options %%%%%
%% The track changes option allows you to mark changes
%% and will produce a list of changes, their line number
%% and page number at the end of the article.
%\documentclass[preprint,trackchanges]{JASAnew}

%% authaffil option will make affil immediately
% follow author, otherwise authors are grouped, and affiliations
% are stacked underneath all the authors.
%\documentclass[preprint,authaffil]{JASAnew}

%% NumberedRefs is used for numbered bibliography and citations.
%% Default is Author-Year style.
% \documentclass[preprint,NumberedRefs]{JASAnew}

%%%%%%% For Reprint
%% For appearance of finished article; 2 columns, 10 pt fonts

% \documentclass[reprint]{JASAnew}

%%%%% Reprint Options %%%%%

%% For testing to see if author has exceeded page length request, use 12pt option
% \documentclass[reprint,12pt]{JASAnew}

% authaffil option will make affil immediately
% follow author, otherwise authors are grouped, and affiliations
% are stacked underneath all the authors.
%\documentclass[reprint,authaffil]{JASAnew}

%% NumberedRefs is used for numbered bibliography and citations.
%% Default is Author-Year style.
\documentclass[reprint,NumberedRefs]{JASAnew}
\usepackage{cases}
%% TurnOnLineNumbers
%% Make lines be numbered in reprint style:
% \documentclass[reprint,TurnOnLineNumbers]{JASAnew}

\begin{document}

\title[JASA/Sample JASA Article]{The Tromba Marina: a physical modelling case study (\#superbadandpretentious)}
\author{Silvin Willemsen}
\email{sil@create.aau.dk}
\affiliation{Multisensory Experience Lab,  Aalborg University Copenhagen, Copenhagen, Denmark}
\author{Michele Ducceschi}
\author{Stefan Bilbao}
\affiliation{Acoustics and Audio Group,  University of Edinburgh, Edinburgh, Scotland}

\author{Stefania Serafin}
\affiliation{Multisensory Experience Lab,  Aalborg University Copenhagen, Copenhagen, Denmark}
 
% \author{Author Five}			% \email{author.five@someplace.edu}
% \affiliation{Department3,  University3, City, State ZipCode, Country}

\preprint{Author, JASA}		%  if you want want this message to appear in upper left corner of title page

\date{\today} 

\begin{abstract}
Put your abstract here. Abstracts are limited to 200 words for
regular articles and 100 words for Letters to the Editor. Please no
personal pronouns, also please do not use the words ``new'' and/or
``novel'' in the abstract. An article usually includes an abstract, a
concise summary of the work covered at length in the main body of the
article.     
\end{abstract}

%% pacs numbers not used

\maketitle

%  End of title page for Preprint option --------------------------------- %


\section{\label{sec:1} Introduction}

\subsection{History}
The tromba marina is a monochord from \textbf{``insert country''} and created in ca. \textbf{``insert year''} (see~\ref{fig:tromba}). It was played mainly by women. The characteristic sound of the instrument comes from the rattling bridge, that makes the instrument sound like a trumpet rather than a bowed-string instrument. 

Non-iterative collisions as presented in \cite{ducceschi2019} \dots

In \cite{bilbao2019}, the same authors have published a continuation with a focus on sonic output quality. To the best of the authors' knowledge, no other literature 


 % Sample showing how to include figure, this is a floating figure
 
\begin{figure}[ht]
%% \reprintcolumnwidth is the same in preprint and reprint for
%% ease of use for authors:
\includegraphics[width=\reprintcolumnwidth]{figsamp.jpg}
\caption{\label{fig:tromba}{The tromba marina.}}
\end{figure}

\section{\label{sec:2} Instrument design and interaction}
The tromba marina is a chordophone... The interaction with the instrument is through bowing. The string rests on a rattling bridge, causing the 
The instrument rests with the neck on the shoulder of the player and the string is bowed close to the neck -- as opposed to other bowed-string instruments which are bowed closer to the bridge. Different pitches are played by playing overtones or harmonics by placing a finger slightly on a one-over-integer multiple of the string length to create a node creating an overtone. \textit{$\Leftarrow$ Notes: need to find better wording.}

The instrument can be divided into three main components: the (bowed) string, the bridge and the body. 

\section{Models}
In this section the model equations will be presented. 
\paragraph{Notation}
Subscripts `$t$' and `$x$' denote a single derivative with respect to time or space. Furthermore, for clarity, subscripts `$\text{s}$' and `$\text{p}$' will be given for parameter symbols that are shared between the string and the plate.
\subsection{Damped Stiff String}
Using state variable $u=u(x,t)$, the motion of a damped stiff string is described as
\begin{equation}
    \rho_\text{s} A u_{tt} = Tu_{xx} - E_\text{s}Iu_{xxxx}-2\rho_\text{s}A\sigma_{0,\text{s}}u_t+2\rho_\text{s}A\sigma_{1,\text{s}}u_{txx},
\end{equation}
with material density $\rho_\text{s}$ (kg$\cdot$m$^{-3}$), cross-sectional area $A = \pi r^2$ (m$^2$), where $r$ is the string radius (m), tension $T$ (N), Young's modulus $E_\text{s}$ (Pa), area moment of inertia $I = \pi r^4 / 4$ (m$^4$) and frequency independent and dependent loss factors $\sigma_{0,\text{s}}$ (s$^{-1}$) and $\sigma_{1,\text{s}}$ (m$^2$/s).
\subsection{Bridge}
The bridge is modelled as a damped mass with a linear restoring force. Its displacement $w=w(t)$ is described as
\begin{equation}
    Mw_{tt}=-M\omega_0^2w-MRw_t,
\end{equation}
with bridge-mass $M$ (kg), angular frequency $\omega_0$ (s$^{-1}$) and damping coefficient $R$ (s$^{-1}$). The restoring force has been added to simulate the fact that\dots  It might seem odd have all terms include a multiplication with mass $M$, but as we start adding the effects of other parts of the system, it will make more sense.
\subsection{Body}
For simplicity, the body is modelled as a 2D plate. The displacement $v = v(x,y,t)$ at location $(x,y)$ is described as
\begin{equation}
    \rho_\text{p} H v_{tt}=-D\Delta\Delta v-2\sigma_{0,\text{p}}v_t+2\sigma_{1,\text{p}}\Delta v_t,
\end{equation}
with material density $\rho_\text{p}$ (kg$\cdot$m$^3$), plate thickness $H$ (m), $D = E_\text{p}H^3/12(1-\nu^2)$ (kg$\cdot$m$^2\cdot$s$^{-2}$), where $E_\text{p}$ is the Young's modulus (Pa) and $\nu$ the unitless Poisson's ratio, and and frequency independent and dependent loss factors $\sigma_{0,\text{p}}$ (s$^{-1}$) and $\sigma_{1,\text{p}}$ (m$^2$/s). Furthermore, $\Delta$ represents the 2D Laplacian and is defined as
\begin{equation}
    \Delta = \frac{\partial^2}{\partial x^2} + \frac{\partial^2}{\partial y^2}.
\end{equation}
\subsection{Implementation}
The complete system is shown in Eqs.~(\ref{eq:systemString}) --~(\ref{eq:systemEtaSpring}).
\begin{subnumcases}{\label{eq:trombaSystem}}
    \begin{aligned}&\rho_\text{s} A \delta_{tt}u^n_l\\
    & \end{aligned} & $\!\!\!\!\!\!\!\!\!\!\begin{aligned}=&\ T\delta_{xx}u^n_l - E_\text{s}I\delta_{xxxx}u^n_l \\
        &- 2\rho_\text{s}A\sigma_{0,\text{s}}\delta_{t\cdot}u^n_l\\
        &+2\rho_\text{s}A\sigma_{1,\text{s}}\delta_{t-}\delta_{xx}u^n_l-J_\text{s}(l_\text{br})F_\alpha\end{aligned}\quad\ \ $\label{eq:systemString} \\
        \begin{aligned}
        &M \delta_{tt}w^n\\
        &
    \end{aligned} &$\!\!\!\!\!\!\!\!\!\!
        \begin{aligned}=&-M\omega_0^2w^n - MR \delta_{t\cdot}w^n\\
            & + \big(\mu_{t+}\psi^{n-1/2}\big)g^n + F_\alpha
        \end{aligned}$\\
    \begin{aligned}
    &\rho_\text{p} H \delta_{tt}v^n_{(l,m)}\\ 
    & \\
    &
    \end{aligned}
    & $\!\!\!\!\!\!\!\!\!\!\begin{aligned} 
    =&-D\delta_{\Delta \boxplus}\delta_{\Delta \boxplus}v_{(l,m)}^n - 2\sigma_{0, \text{p}}\delta_{t\cdot}v^n_{(l,m)}\\
    &%\!\!\!\!\!\!\!\!\!\!\!\!\!\!\!\!\!\!\!\!\!\!\!\!\!\!
    + 2\sigma_{1, \text{p}}\delta_{t-}\delta_{\Delta\boxplus}v_(l,m)^n\\
    & - J_\text{p}(l_\text{br},m_\text{br}) \big(\mu_{t+}\psi^{n-1/2}\big)g^n
    \end{aligned}$\\
\begin{aligned}
\delta_{t+}\psi^{n-1/2} 
\end{aligned}
&$\!\!\!\!\!\!\!\!\!\!
\begin{aligned}= g^n\delta_{t\cdot}\eta_\text{c}^n
\end{aligned}$\\
\begin{aligned}
    \eta_c^n 
\end{aligned}
&$\!\!\!\!\!\!\!\!\!\!
\begin{aligned}
    = v_{(l_\text{br},m_\text{br})}^n - w^n
\end{aligned}$\\
    \begin{aligned}
        &F_\alpha \\
        &
    \end{aligned}
    & $\!\!\!\!\!\!\!\!\!\!
    \begin{aligned}
        =&\ K_1\mu_{tt}\eta_\text{sp}^n + K_3(\eta_\text{sp}^n)^2\mu_{t\cdot}\eta_\text{sp}^n \\
        &+ 2 \sigma_\times\delta_{t\cdot}\eta_\text{sp}^n
    \end{aligned}$\label{eq:fAlphaDamp}\\
\begin{aligned}\eta_s^n\end{aligned} &$\!\!\!\!\!\!\!\!\!\!
\begin{aligned}
    = u_{l_\text{br}}^n - w^n
\end{aligned}$\label{eq:systemEtaSpring}
\end{subnumcases}
% \begin{subnumcases}{\label{eq:trombaSystem}}
%     \rho_\text{s} A \delta_{tt}u^n_l & $\!\!\!\!=T\delta_{xx}u^n_l - EI\delta_{xxxx}u^n_l - J_\text{s}(l_\text{br})F_\alpha\qquad\quad$ \\
%     M \delta_{tt}w^n &$\!\!\!\!=-M\omega_0^2w^n + \big(\mu_{t+}\psi^{n-1/2}\big)g^n + F_\alpha\qquad\quad$\\
%     \begin{aligned}&\rho_\text{p} H \delta_{tt}v^n_{(l,m)}\\ 
%     & \end{aligned} & $\!\!\!\!\!\!\!\,\,\begin{aligned} =&-D\delta_{\Delta \boxplus}\delta_{\Delta \boxplus}v_{(l,m)}^n \\
%     &- J_\text{p}(l_\text{br},m_\text{br}) \big(\mu_{t+}\psi^{n-1/2}\big)g^n \end{aligned}$\\
%     \delta_{t+}\psi^{n-1/2} &$\!\!\!\!= g^n\delta_{t\cdot}\eta_\text{c}^n$\\
%     \eta_c^n &$\!\!\!\!= v_{(l_\text{br},m_\text{br})}^n - w^n$\\
%     F_\alpha & $\!\!\!\!= K_1\mu_{tt}\eta_\text{s}^n + K_3(\eta_\text{s}^n)^2\mu_{t\cdot}\eta_\text{s}^n$\label{eq:fAlpha}\\
%     \eta_s^n &$\!\!\!\!= u_{l_\text{br}}^n - w^n$\label{eq:etaSpring}
% \end{subnumcases}
\\
%
The time-difference operators have been chosen so that the system is most accurate but still explicit 

\section{Floats, Figures and Tables}

 Figures and tables are typically ``floats'' which means that their
final position is determined by \LaTeX\ while the document is being
typeset. \LaTeX\ isn't always successful in placing floats
optimally.  Use the figure* environment to get a wide figure that spans the 
page in a two-column layout.

\subsection{\label{subsec:3:2} Tables}
Tables generally should be surrounded with
\verb+\begin{ruledtabular}...\end{ruledtabular}+\\
This will guarantee that they are the width of the
page or column, and have two ruled lines at the top
and bottom of the table.

\verb+[ht]+ in the code below instructs \LaTeX\ to place the table
where it appears in type, if it will fit on the page;
otherwise put it on the top of the next page.

Footnotes in a table are labeled a, b, c,
 etc.  They can
be  specified  by  using  the  \LaTeX\
\verb+\footnotemark[]+
and
\verb+\footnotetext[]+ commands.
The  footnotes  for  a  table  are  typeset  at  the
bottom  of  the  table,  rather  than  at  the  bottom  of  the
page or at the end of the references.  The arguments for
\verb+\footnotemark[]+
and
\verb+\footnotetext[]+
should be numbers
1, 2, \ldots\  The journal style will convert these to letters.

This system allows
multiple entries to refer to the same
footnote.   


\begin{table}[ht]
\caption{\label{tab:table1}A table with more columns still fits
properly in a column. Note that several entries share the same
footnote. Inspect the \LaTeX\ input for this table to see
exactly how it is done.}

\begin{ruledtabular}
\begin{tabular}{cccccccc}
 &$r_c$ (\AA)\footnotemark[1]&$r_0$ (\AA)&$\kappa r_0$&
 &$r_c$ (\AA) &$r_0$ (\AA)&$\kappa r_0$\\
\hline
Cu& 0.800 & 14.10 & 2.550 &Sn\footnotemark[1]
& 0.680 & 1.870 & 3.700 \\
Ag& 0.990 & 15.90 & 2.710 &Pb\footnotemark[2]
& 0.450 & 1.930 & 3.760 \\
%Au& 1.150 & 15.90 & 2.710 &Ca\footnotemark[3]
%& 0.750 & 2.170 & 3.560 
\end{tabular}
\end{ruledtabular}
\footnotetext[1]{Here's the first.}
\footnotetext[2]{Here's the second.}
%\footnotetext[3]{Here's the third.}
\end{table}
\clearpage
\subsection{Plain Tables: When NOT to use `ruledtabular'}
There are a number of cases when `ruledtablar' should not
be used: basically for any table using complex content or 
commands.

\subsubsection{Using {\tt\string\multicolumn}}
When you'd like to use the multicolumn command in your
table, you'll find that `ruledtabular' will cause bad
formatting. In that case, Don't Use Ruledtabular, and
instead put in \verb+\hline\hline+ at the top and bottom of
the table. 


\begin{table}[ht]
\caption{A table made without `ruledtabular' needs to have two hlines
added to the top and bottom of the table.}
\vskip3pt
\begin{tabular}{cccccccc}
\hline\hline
 &$r_c$ (\AA)\footnotemark[1]&$r_0$ (\AA)&$\kappa r_0$&
 &$r_c$ (\AA) &$r_0$ (\AA)&$\kappa r_0$\\
\hline
Cu& 0.800 & 14.10 & 2.550 &Sn\footnotemark[1]
& 0.680 & 1.870 & 3.700 \\
Ag& 0.990 & 15.90 & 2.710 &Pb\footnotemark[2]
& 0.450 & 1.930 & 3.760 \\
Au& 1.150 & 15.90 & 2.710 &Ca\footnotemark[3]
& 0.750 & 2.170 & 3.560\\
\hline\hline
\end{tabular}
\footnotetext[1]{This is the first table note.}
\footnotetext[2]{This is the second table note.}
\footnotetext[3]{This is the third table note.}
\end{table}

\newpage

\vskip12pt
\vskip1sp
\subsubsection{Using the {\tt\string\adjustbox\string{\string}
(tabular)\string\end\string{adjustbox\string}}
command}
There may be times when the table is too wide, or you want to have
the table be the width of the page, whether or not it appears in
preprint or reprint version of JASAnew. In this case you can use
\verb+\begin{adjustbox}{<width>} (tabular) \end{adjustbox}+.
(`adjustbox' will NOT work with `ruledtabular')

You can set a maximum width with \\
\verb+\begin{adjustbox}{max width=\textwidth}(tabular)\end{adjustbox}+\\  
in which case the table in the reprint version might be less than
the full text width;

Or you can set the exact width you'd like with\\
\verb+\begin{adjustbox}{width=\textwidth}(tabular)\end{adjustbox}+\\
in which case the table will be the full width of the page in
either preprint or reprint.

This way you can make a table that will fit in the correct width
whether you are using the preprint or reprint option.

\begin{table*}[b]
\caption{Top~5 rated $\widehat{\textrm{ITD}}$ estimation methods
according to the sum and product metric criteria for $\pm0.5$~JND and
$\pm1$~JND tolerance thresholds (normalized scores).}

\begin{adjustbox}{width=\textwidth}
\begin{tabular}{c| l l| l l| l l| l l}
\hline\hline
Rank \#& \multicolumn{2}{c}{sum criteria [$\pm0.5$~JND]} &
\multicolumn{2}{|c}{sum criteria [$\pm1$~JND]} &
\multicolumn{2}{|c}{product criteria [$\pm0.5$~JND]} & \multicolumn{2}{|c}{product criteria [$\pm1$~JND]}\\ 
\hline
1 & Threshold --30dB {lp} & (0.43) & Threshold --30dB {lp} & (0.71) & Threshold --30dB {lp} & (1.00) & Threshold --30dB {lp} & (1.00) \\ 
2 & MaxIACCe {lp} & (0.39) & Threshold --20dB {lp} & (0.66) & MaxIACCe {lp} & (0.39) & Threshold --20dB {lp} & (0.57) \\ 
3 & Threshold --20dB {lp} & (0.38) & CenIACCr {bb} &(0.62) & CenIACCr {lp} &(0.33) & CenIACCr {bb} & (0.37) \\ 
4 & CenIACCr {lp} &(0.37) & MaxIACCe {lp} & (0.61) & Threshold --20dB {lp} & (0.29) & MaxIACCe {lp} & (0.34) \\ 
5 & Cen-$e^2$ {lp} &(0.34) & CenIACCe {lp} & (0.61) & Cen-$e^2$ {lp}
& (0.10) & CenIACCr {lp} & (0.33) \\
\hline\hline
\end{tabular}
\end{adjustbox}

\label{tab:topMethods} 
\end{table*}



\clearpage


\subsection{Using dcolumn}
\verb+\usepackage{dcolumn}+ is included in JASAnew.cls so you
don't need to add it explicitly.
\href{http://anorien.csc.warwick.ac.uk/mirrors/CTAN/macros/latex/required/tools/dcolumn.pdf}
{\url{http://anorien.csc.warwick.ac.uk/mirrors/CTAN/macros/latex/required/tools/dcolumn.pdf}}
will give you detailed information.
A gentler introduction may be found in this informative and well
illustrated
article: \href{https://www.tug.org/pracjourn/2007-1/mori/mori.pdf}
{\url{https://www.tug.org/pracjourn/2007-1/mori/mori.pdf}}, starting on page
20. (You may want to look at more examples in this quite comprehensive
article on making tables in \LaTeX.)

\begin{quote}
``If we do not want to break the fractional and the integral part in two columns,
the dcolumn package provides a new type of column\\
\verb+D{sep -in}{sep -out}{ before.after}+\\
The first argument \verb+{sep-in}+ is the symbol used in the
.tex document to separate
the integral and the fractional part (usually the decimal point . or the decimal
comma ,), the second argument \verb+{sep-out}+
is the symbol that we want in the
output, the third is the number of digits on the left (before) and on the right
(after) this symbol. The numbers are aligned to the decimal point and, in case
that the third argument is negative, the decimal point is aligned to the center of
the column. If the columns have a heading, it must be inserted into
the command \verb+\multicolumn{1}{c}{...}+
\end{quote}
\newpage
An example using dcolumn:
\begin{verbatim}
{\hsize= 2in
\begin{ruledtabular}
\begin{tabular}{cD {,}{.}{5.4}}
Expression           & \multicolumn {1}{c}{ Value }\\
\hline
$\pi$                  &      3,1416                 \\
$\pi^{\pi}$           &     36,46                    \\
$\pi^{\pi^{\pi}}$    & 80662,7                      \\
\end{tabular}
\end{ruledtabular}
}
\end{verbatim}
\vskip12pt
{\hsize= 2in
\begin{ruledtabular}
\begin{tabular}{cD {,}{.}{5.4}}
Expression           & \multicolumn {1}{c}{ Value }\\
\hline
$\pi$                  &      3,1416                 \\
$\pi^{\pi}$           &     36,46                    \\
$\pi^{\pi^{\pi}}$    & 80662,7                      \\
\end{tabular}
\end{ruledtabular}
}



\newpage



 
\begin{figure*}[p]
\parskip=6pt
\baselineskip=12pt



\subsection{Sample Figures, new commands available in this style}
\vglue-14pt
{\it Note that the publisher determines the final layout, so 
your choice of figure alignment may not be reflected in the published
article.}

\noindent
\verb+\figline{}+ will center one or more figures on one line. 


\noindent
\verb+\fig{<name of file>}{<width>}{<letter to put underneath>}+

\noindent
\verb+\leftfig{<name of file>}{<width>}{<letter to put underneath>}+

\noindent
\verb+\rightfig{<name of file>}{<width>}{<letter to put underneath>}+

\noindent
\verb+\boxedfig{<name of file>}{<width>}{<letter to put underneath>}+

\noindent
\verb+\rotatefig{<degrees of rotation>}{<name of file>}{<width>}+\\
\verb+{<letter to put underneath>}+

The following illustrations show these commands in use.


\begin{verbatim}
\figline{\fig{figsamp.jpg}{4cm}{(a)}
\fig{figsamp.jpg}{4cm}{(b)}}
\figline{\fig{figsamp.jpg}{4cm}{(c)}
\fig{figsamp.jpg}{4cm}{(d)}}
\figline{\fig{figsamp.jpg}{4cm}{(e)}}
\end{verbatim}

\figline{\fig{figsamp.jpg}{4cm}{(a)}
\fig{figsamp.jpg}{4cm}{(b)}}
\figline{\fig{figsamp.jpg}{4cm}{(c)}
\fig{figsamp.jpg}{4cm}{(d)}}
\figline{\fig{figsamp.jpg}{4cm}{(e)}}

\caption{ \label{fig:pressure_field} Multiple images on one figure example (a) image 1, (b-f) ($\rho$=1000 kg/m$^3$) and speed of sound ($c$=1500 m/s).}

\end{figure*}

\begin{figure*}
\baselineskip=12pt
\begin{verbatim}
\figline{\boxedfig{figsamp.jpg}{2in}{(a)}}
\figline{\leftfig{figsamp.jpg}{2in}{(b)}\rightfig{figsamp.jpg}{2in}{(c)}}
\figline{\rotatefig{90}{figsamp.jpg}{2in}{(d)}\rotatefig{180}{figsamp.jpg}{2in}{(e)}}
\end{verbatim}

\figline{\boxedfig{figsamp.jpg}{2in}{(a)}}

\figline{\leftfig{figsamp.jpg}{2in}{(b)}\rightfig{figsamp.jpg}{2in}{(c)}}

\figline{\rotatefig{90}{figsamp.jpg}{2in}{(d)}
\rotatefig{180}{figsamp.jpg}{2in}{(e)}
}

\caption{More figure examples: (a) boxedfig, 
(b)leftfig; (c)right fig; (d) rotatefig 90 degrees;
(e) rotatefig 180 degrees. }

\end{figure*}








\begin{figure*}
\baselineskip=12pt
\begin{verbatim}
\sidebysidefigures{figsamp.jpg}{Describing the first
illustration.}/{figsamp.jpg}{Describing the second illustration.}
\end{verbatim}

\sidebysidefigures{figsamp.jpg}{Describing the first
illustration.}/{figsamp.jpg}{Describing the second illustration.}
\end{figure*}


\begin{figure*}
\baselineskip=12pt
\begin{verbatim}
\figline{
\fig{figsamp.jpg}{.7\textwidth}{}
\narrowcaption{.2\textwidth}{Here is a narrow caption.}
}
\end{verbatim}
\figline{\fig{figsamp.jpg}{.7\textwidth}{} \narrowcaption{.2\textwidth}{Here is a narrow
caption.}}
\end{figure*}


\begin{figure*}
\baselineskip=12pt
\begin{verbatim}
\figline{\fig{figsamp.jpg}{.2\textwidth}{(A)}
\fig{figsamp.jpg}{.2\textwidth}{(B)}
\fig{figsamp.jpg}{.2\textwidth}{(C)}
\narrowcaption{.25\textwidth}{Caption for three illustrations. 
The caption may produce many lines, but only one paragraph.
}}
\end{verbatim}
\figline{\fig{figsamp.jpg}{.2\textwidth}{(A)}
\fig{figsamp.jpg}{.2\textwidth}{(B)}
\fig{figsamp.jpg}{.2\textwidth}{(C)}
\narrowcaption{.25\textwidth}{Caption for three illustrations. 
The caption may produce many lines, but only one paragraph.
}}
\end{figure*}



\begin{figure*}
\baselineskip=12pt
\begin{verbatim}
\figline{\fig{figsamp.jpg}{.7\textwidth}{}
\narrowcaption{.25\textwidth}{Here is a narrow caption that will can be
positioned to the right of four illustrations.
You cannot have more than one paragraph of text in a caption.
You cannot have more than one paragraph of text in a caption.
You cannot have more than one paragraph of text in a caption.
You cannot have more than one paragraph of text in a caption.
}}
\end{verbatim}
\figline{\fig{figsamp.jpg}{.7\textwidth}{}
\narrowcaption{.25\textwidth}{Here is a narrow caption that will can be
positioned to the right of four illustrations.
You cannot have more than one paragraph of text in a caption.
You cannot have more than one paragraph of text in a caption.
You cannot have more than one paragraph of text in a caption.
You cannot have more than one paragraph of text in a caption.
}}
\end{figure*}
\clearpage

\begin{figure}[h]
\baselineskip=12pt
\begin{verbatim}
\figcolumn{
\fig{figsamp.jpg}{.2\textwidth}{(A)}
\fig{figsamp.jpg}{.2\textwidth}{(B)}
\fig{figsamp.jpg}{.2\textwidth}{(C)}
}
\end{verbatim}

\figcolumn{
\fig{figsamp.jpg}{.2\textwidth}{(A)}
\fig{figsamp.jpg}{.2\textwidth}{(B)}
\fig{figsamp.jpg}{.2\textwidth}{(C)}
}

\caption{Here are some stacking figures in a single column. 
The {\tt\string\figcolumn\string{\string}} command works equally well in single
or double column figures.}
\end{figure}

\subsection{Example of multimedia entry}
Please note that this is for multimedia intended to appear inline
within the published article. 

Here is what a multimedia entry will look like:
\multimedia{http://dx.doi.org/10.1121/1.4947423.1}{Corresponding pulse-compressed echo envelopes
and video recordings from a fluttering luna moth.
Echoes from the wings and body of the moth generally dominate the
acoustic returns, which vary greatly over consecutive ensonifications
across the wingbeat cycle. File of type ``mp4'' (15.3
MB)}\label{mmtest1}

Here we try cross referencing the multimedia entry: The multimedia
above is Mm.~\ref{mmtest1}.

\subsection{Supplementary Material}
ASA
prefers that authors to submit related/relevant article files as
supplementary material with their submission.

\subsection{Supplementary material for publication}
Any archival supplemental materials to be published with the
manuscript (eg., supplementary figures) should be cited in-text and a footnote provided.

An example of reference to supplementary material:

The sound files and videos for this and other figures
are included as supplementary materials\footnote{See
Supplementary materials at [URL will be inserted by AIP]
for [give a brief description of the material].}.

The contents of the footnote above will appear at the beginning of the
bibliography made with BibTeX when the default `author-year' documentclass option is used;
BibTeX output will have the footnote interleaved with other
references if the NumberedRefs documentclass option is used.

\subsection{File naming conventions}
Here are the conventions for naming files:

\begin{itemize}
\item
Supplementary Figure or
	Supplementary Figure or Text files should be named: SuppPub\#.xxx, where ``\#'' is
	a number and ``xxx" is the file format extension
	(SuppPub1.docx, SuppPub2.jpg, etc)

\item
	Supplementary Multimedia files: SuppPubmm\#.xxx, where ``\#'' is a
	number and ``xxx'' is the file format extension (SuppPubmm1.mp3,
	SuppPubmm2.gif, etc)

\item
Multimedia files must be named accordingly: MM\#.xxx, where ``\#'' is the
number and ``xxx'' is the file format extension (MM1.wav, MM2.avi, etc).

\item
The only figure formats allowed are the following: 
.pdf, .ps, .eps, or .jpg. Figure files must be named in this fashion:
Figure\#.xxx, where ``\#'' is the figure number and ``xxx'' is the file format
(Figure1.eps, Figure2.jpg, Figure3a.ps, Figure3b.ps, etc). 

\end{itemize}



\section{\label{sec:5}Conclusion}

And in conclusion\ldots

\begin{acknowledgments}
This research was supported by  ...
\end{acknowledgments}


 % -------------------------------------------------------------------------------------------------------------------
 %   Appendix  (optional)


\appendix
\section{Appendices}
To start the appendix, use the \verb+\appendix+ command.
This signals that all following section commands refer to appendixes
instead of regular sections. Therefore, the \verb+\appendix+ command
should be used only once---to set up the section commands to act as
appendices. Thereafter normal section commands are used. The heading
for a section can be left empty. For example,
\begin{verbatim}
\appendix
\section{}
\end{verbatim}
will produce an appendix heading that says ``APPENDIX A'' and
\begin{verbatim}
\appendix
\section{Background}
\end{verbatim}
will produce an appendix heading that says ``APPENDIX A: BACKGROUND''
(note that the colon is set automatically).

If there is only one appendix, then the letter ``A'' should not
appear. This is suppressed by using the star version of the appendix
command (\verb+\appendix*+ in the place of \verb+\appendix+).
\vskip12pt
\vskip1sp
\section{A little more on appendices}

Observe that this appendix was started by using
\begin{verbatim}
\section{A little more on appendixes}
\end{verbatim}

Note the equation number in an appendix:
\begin{equation}
E=mc^2.
\end{equation}


\subsection{\label{app:subsec}A subsection in an appendix}

You can use a subsection or subsubsection in an appendix. Note the
numbering: we are now in Appendix~\ref{app:subsec}.

\subsubsection{\label{app:subsubsec}A subsubsection in an appendix}
Note the equation numbers in this appendix, produced with the
subequations environment:
\begin{subequations}
\begin{eqnarray}
E&=&mc, \label{appa}
\\
E&=&mc^2, \label{appb}
\\
E&\agt& mc^3. \label{appc}
\end{eqnarray}
\end{subequations}
They turn out to be Eqs.~(\ref{appa}), (\ref{appb}), and (\ref{appc}).


\section{Sample appendix figure and table}
Figure and table numbering are continuous through the article,
and handled the same as they are in the rest of the article.


\begin{figure}[ht]
\includegraphics[width=1in]{figsamp.jpg}
\caption{Figure in an appendix.}
\end{figure}

\begin{table}[ht]
\caption{Here is the caption for a table in an appendix.}
\centering
\begin{ruledtabular}
\begin{tabular}{cccc}
one&two&three&four\\
\hline
C&D&E&F\\
\end{tabular}
\end{ruledtabular}
\vskip12pt
\end{table}

\subsection{Footnotes}
The contents of the footnotes will appear at the beginning of the
bibliography when BibTeX produces the .bbl file using the default
AuthorYear style; interleaved with other references if 
Numbered\-Refs option:\\
\verb+\documentclass[preprint,NumberedRefs]{JASAnew}+\\
and BibTeX has been used.

This example show where this cite \citep{booksamp1} will appear in the
bibliography,
depending on whether we use default author-year style
or call for the NumberedRefs documentclass option.

Here are some sample footnotes:\footnote{Here is the second footnote.
It will appear before the beginning of the bibliography in Author-Year
style (default) or it will be 
 interleaved with other references when using the Numbered\-Refs
 option.}$^,$% (comma between two sequential footnotes)
\footnote{Here is a third footnote.}


\section{Making the Bibliography Using BibTeX}
Authors are highly  recommended to use BibTeX to produce their
bibliographies. The results will be predictable and even if
it might take some time to get comfortable with  using BibTeX,
in the long run it will save you endless aggravation.

A resource for making your bibliography entries
correctly is included in this package: 
JASA-ReferenceStyles.pdf. You will also find
the files
bibsamp1.tex/.pdf and bibsamp2.tex/.pdf
for examples of output; and sampbib.bib for an example of
how to make your .bib database entries.

There are two possible bibliography styles: the default, author-year,
and the optional style, Numbered\-Refs, which you would call using\\
{\verb+\documentclass[preprint,NumberedRefs]{JASAnew}+ }

Every \verb+\cite+ will produce a citation and an entry in the
bibliography and every cite must have a matching entry in the bib
database file.

\verb+\citep{}+ should be used rather than \verb+\cite{}+
Note that the citations are hyperlinked to their entries in the
bibliography:

Normal journal cite: \citep{joursamp1},
 Book reference \citep{booksamp1},
In press, \citep{inpress3}. 
Computer language documentation:
\citep{sampcode2}.


\begin{quote}
\it
NOTE:\\
Once you have used BibTeX you
should open the resulting .bbl file and cut and paste the entire contents 
into the end of your article. You should also comment out
\verb+\bibliography{<your .bib file>}+, ie,
\verb+%\bibliography{<your .bib file>}+.
\end{quote}

Make your bibliography by doing: pdflatex filename,  bibtex filename,
pdflatex filename, pdflatex filename.

{\bfseries\itshape
Compare the results you get with\\
{\verb+\documentclass[preprint]{JASAnew}+ }\\
vs.\\
{\verb+\documentclass[preprint,NumberedRefs]{JASAnew}+ }
}

\bibliography{sampbib}


\end{document}




